\chapter{Trazabilidad de requisitos y evidencias}
\label{chap:anexo-trazabilidad}

\section{Propósito del anexo}
\label{sec:anexo-proposito-trazabilidad}

Este anexo resume la relación entre requisitos, implementación, validación y evidencia. Su objetivo es facilitar la defensa del proyecto y evitar que la evaluación dependa únicamente de una explicación oral. Cada requisito importante puede relacionarse con una parte del código, un mecanismo de prueba y una evidencia documental o cloud.

La trazabilidad no sustituye a los documentos completos del repositorio. La matriz operativa principal se mantiene en \path{docs/RequirementsTraceability.md}; este anexo recoge una versión integrada en la memoria para que el PDF sea autocontenido.

\section{Requisitos funcionales}
\label{sec:anexo-rf}

{\setlength{\tabcolsep}{4pt}\footnotesize
\begin{longtable}{@{}p{0.08\textwidth}p{0.28\textwidth}p{0.27\textwidth}p{0.25\textwidth}@{}}
\caption{Trazabilidad de requisitos funcionales}
\label{tab:trazabilidad-rf}\\
\textbf{ID} & \textbf{Implementación} & \textbf{Validación} & \textbf{Evidencia} \\
\hline
RF-01 & Endpoints de salud \path{/api/v1/health/live} y \path{/api/v1/health/ready}. & Tests de health, smoke local y smoke cloud. & Health HTTPS \texttt{200}, target group sano y logs de readiness. \\
RF-02 & Endpoint de bootstrap protegido por token externo. & Pruebas de creación inicial y conflicto si ya existe administrador. & Bootstrap cloud \texttt{201} o \texttt{409} documentado. \\
RF-03 & Login con usuario y contraseña, emisión de JWT. & Tests de autenticación y smoke cloud de login. & Login administrador \texttt{200} con JWT. \\
RF-04 & Autorización por roles \texttt{admin} y \texttt{operator}. & Tests de endpoints protegidos y acceso denegado. & Casos de prueba y documentación de roles. \\
RF-05 & Endpoints de administración de usuarios. & Tests de creación, listado, cambio de rol y activación. & Suite automatizada de backend. \\
RF-06 & Gestión de transportes con filtros, paginación y borrado lógico. & Tests de CRUD, validación y soft delete. & Pruebas de integración. \\
RF-07 & Gestión de vehículos asignables. & Tests de validación de matrícula, unicidad y borrado lógico. & Pruebas de dominio y persistencia. \\
RF-08 & Gestión de conductores asignables. & Tests de licencia, contacto, estado activo y unicidad. & Pruebas de integración. \\
RF-09 & Asignación de vehículo y conductor a transporte. & Tests de asignación válida e inválida. & Validación de reglas de negocio. \\
RF-10 & Transiciones de estado controladas. & Tests de estados terminales y precondiciones. & Suite de reglas de transporte. \\
RF-11 & Registro y consulta de eventos logísticos. & Tests de creación y consulta cronológica. & Historial de eventos validado. \\
RF-12 & Errores coherentes y contrato común. & Tests de validación, no encontrado, conflicto y autorización. & Respuestas HTTP y request id correlable. \\
\end{longtable}
}

\section{Requisitos no funcionales}
\label{sec:anexo-rnf}

{\setlength{\tabcolsep}{4pt}\footnotesize
\begin{longtable}{@{}p{0.14\textwidth}p{0.29\textwidth}p{0.24\textwidth}p{0.21\textwidth}@{}}
\caption{Trazabilidad de requisitos no funcionales}
\label{tab:trazabilidad-rnf}\\
\textbf{Área} & \textbf{Implementación} & \textbf{Validación} & \textbf{Evidencia} \\
\hline
Reprod. local & Docker Compose con API y PostgreSQL. & Smoke local y pruebas automatizadas. & Comandos documentados en README y validación local. \\
Mantenibilidad & Solución simple con proyectos \path{TransitOps.Api} y \path{TransitOps.Tests}. & Revisión de estructura y tests. & Capítulos de arquitectura e implementación. \\
Persistencia & PostgreSQL y migraciones EF Core. & Migraciones locales y ECS \path{--migrate-only}. & Tareas ECS de migración con exit code \texttt{0}. \\
Seguridad básica & Hash de contraseñas, JWT, roles, secretos externos y red privada. & Tests de auth, auditoría Sprint 8 y revisión de SG/IAM. & Script de posture audit y documentación de seguridad. \\
IaC & Terraform con módulos y estado remoto. & \texttt{terraform fmt}, validate, apply y destroy. & Recreate-from-scratch y state vacío tras destroy. \\
Observab. & Logs JSON, correlation id, dashboard, alarmas y filtro métrico. & Smoke con correlation id y verificación CloudWatch. & Dashboard, 9 alarmas y logs filtrables. \\
CI/CD & Workflows manuales de Terraform, deploy y rollback. & Ejecución local, publicación de imagen y validación cloud. & Documentación de pipelines y capturas de ECR, ECS, readiness y task definition activa. \\
Operación & Runbooks de deploy, rollback, restore, logs, alarmas y destroy. & Sprint 7 y Sprint 8 con AWS real. & \path{docs/CloudReliability.md}, \path{docs/CloudDeployment.md}. \\
Coste & Entorno \texttt{dev} efímero, backups RDS a cero y destroy final. & Auditoría post-destroy. & Ausencia de recursos efímeros costosos. \\
\end{longtable}
}

\section{Evidencias por sprint}
\label{sec:anexo-evidencias-sprint}

{\setlength{\tabcolsep}{4pt}\footnotesize
\begin{longtable}{@{}p{0.11\textwidth}p{0.32\textwidth}p{0.45\textwidth}@{}}
\caption{Resumen de evidencias por sprint}
\label{tab:evidencias-sprint}\\
\textbf{Sprint} & \textbf{Objetivo} & \textbf{Evidencia generada} \\
\hline
Sprint 1 & MVP local backend. & Endpoints funcionales, modelo de dominio, pruebas iniciales y ejecución local. \\
Sprint 2 & Base Terraform y red AWS. & Estructura de módulos, remote state, VPC, subredes, NAT y grupos de seguridad. \\
Sprint 3 & Runtime cloud planificado. & ECR, RDS, ECS, ALB, Route 53, ACM, Secrets Manager y CloudWatch definidos. \\
Sprint 4 & Primer despliegue real. & Imagen en ECR, migraciones ECS, servicio estable, HTTPS y login admin. \\
Sprint 5 & CI/CD con OIDC. & Workflows de validación y despliegue manual sin claves AWS persistentes. \\
Sprint 6 & Observabilidad y hardening. & Logging JSON, correlation id, dashboard, alarmas, circuit breaker y DNS/HTTPS corregido. \\
Sprint 7 & Rollback, restore y tuning. & Rollback a tag bueno, fallo con tag inexistente, restore RDS temporal y Docker hardening. \\
Sprint 8 & Seguridad, coste y runbooks. & Auditoría de postura, recreate-from-scratch, coste documentado y destroy audit. \\
Sprint 9 & Cierre documental. & Trazabilidad, índice final de evidencias, guion de demo y memoria final. \\
\end{longtable}
}

\section{Relación entre evidencia y defensa}
\label{sec:anexo-defensa}

Las evidencias se agrupan en tres niveles. El primer nivel corresponde a evidencias locales: tests automatizados, build Docker, usuario no root y validación de Terraform. Estas pruebas demuestran que el repositorio es coherente y que el sistema puede construirse sin depender de AWS.

El segundo nivel corresponde a evidencias cloud: Terraform apply, imagen ECR, migración ECS, servicio estable, health HTTPS, login, logs correlados, dashboard, alarmas, rollback y restore. Estas evidencias demuestran que el proyecto no se queda en diseño teórico y que se ejecutó sobre servicios reales.

El tercer nivel corresponde a evidencias de cierre: destroy final, state vacío, ausencia de recursos costosos, coste revisado y recursos persistentes justificados. Este nivel es importante porque una arquitectura cloud académica no debe ignorar el gasto ni dejar recursos activos por descuido.

Durante la defensa, estas evidencias permiten recorrer el proyecto desde requisitos hasta operación. Un posible guion consiste en explicar primero el problema y el alcance, después enseñar la arquitectura y el flujo de despliegue, luego validar una petición con correlation id y finalmente mostrar rollback, restore y destroy como pruebas de madurez operativa.

\section{Capturas incorporadas}
\label{sec:anexo-capturas-incorporadas}

La memoria incorpora capturas reales de las validaciones cloud y de las auditorías finales. Las evidencias principales incluyen:

\begin{itemize}
  \item Servicio ECS estable con una tarea en ejecución.
  \item Certificado ACM en estado \texttt{Issued}.
  \item Route 53 con el registro del subdominio y validación ACM.
  \item Dashboard CloudWatch.
  \item Consulta de logs por \httpheader{X-Correlation-ID}.
  \item Publicación en ECR, task definition activa y readiness HTTPS.
  \item Prueba de restore RDS.
  \item Auditoría de seguridad Sprint 8.
  \item Cost Explorer durante recreación temporal.
  \item Auditoría post-destroy.
\end{itemize}

Estas capturas permiten defender el recorrido completo: despliegue, operación, observabilidad, recuperación, coste y cierre del entorno.
